\documentclass[10pt,a4paper]{amsart}
\usepackage[margin=19mm]{geometry}
\usepackage{amsmath,amssymb,mathtools}
\usepackage[scr=rsfs,cal=euler]{mathalfa}
\usepackage{xfrac}
\usepackage[unicode,hidelinks,bookmarks=false]{hyperref}
\newcommand{\ie}{i.e.\ }
\newcommand{\cf}{cf.\ }
\newcommand{\resp}{respectively\ }
\newcommand{\Subsection}{\subsection}
\providecommand{\nobreakdash}{\nobreak\hbox{-}}
\setlength{\emergencystretch}{2em}
\multlinegap0pt
\usepackage{amssymb}
\usepackage{bm}
\usepackage{xcolor}
\usepackage{algorithm}\usepackage{algpseudocode}
\newtheorem{problem}{Problem}
\newtheorem{theorem}{Theorem}
\def\w{\widetilde}
\title{Generalized Fine-Tuning of Diffusion Models via Stochastic Control and FBSDEs}
\author{Zirui Wang, Lu Wang}
\date{Source: arXiv:2607.22660v1 (CC BY 4.0)}
\begin{document}
\maketitle

\noindent\emph{Source and licence.} Generalized Fine-Tuning of Diffusion Models via Stochastic Control and FBSDEs. Version arXiv:2607.22660v1 (CC BY 4.0), distributed by arXiv under the Creative Commons Attribution 4.0 International licence, http://creativecommons.org/licenses/by/4.0/, as recorded in the arXiv licence metadata for this exact version and shown on the abstract page https://arxiv.org/abs/2607.22660v1. Full licence notice: \texttt{licenses/arxiv-2607.22660-CC-BY-4.0.txt}.\par\smallskip

\noindent\textit{Editorial scope.} Counted unit: the article's theorem theo:Regularity (source0.tex lines 486--521). The article's complete original proof of it is delivered with it (source0.tex lines 524--667, one \texttt{proof} environment of 144 lines). No mathematics is written, repaired or re-derived: the statement and the proof are the article's own lines, in the article's own notation, and no retained line is substituted or re-flowed.  Retained uncounted as the source-local context the proof needs: lines 189--193; lines 444--450; lines 458--461; lines 463--466.  Nothing was omitted from any retained span, so the package carries no omission record.  No retained line contains a citation, so the wrapper carries no bibliography and no bibliography entry.  No retained line contains a figure environment, an image-inclusion command or drawing code, so no image is delivered and the package declares no asset dependency.  Editorial formatting only, in addition: an AMS \texttt{amsart} preamble that restates the article's own declarations and macros used by the retained lines, namely the environments \texttt{problem}, \texttt{theorem} on separate counters, together with the standard packages those lines need.  The retained environments print in this document as Problem 1, Theorem 1 in order of appearance, whereas the article prints the same items with the numbering of its own full document.  The complete native source of the article as distributed in the arXiv e-print for this exact version is bundled unchanged as \texttt{sources/205966/source0.tex}.
\begin{equation}
\label{state}
d X_t = [\w{b}(t,{X}_t)+u_t ]\, dt +  \w{\sigma}(t), dB_t
\quad X_0 \sim  p_{\text{noise}}(\cdot)
\end{equation}

\begin{problem}
\textbf{GFT} 
\label{prob:GFT}

Find an optimal $u^*$ such that 
$\widetilde{J}(u^*) = \inf_{u}\widetilde{J}(u)$ subject to the state dynamics given in \eqref{state}.    
\end{problem}

\begin{equation}
\label{eq:v^*}
v^*(t,x) := \inf_{u} \mathbb{E}\left[ \int_t^T \ell(s, X_s, u_s)ds + r(X_T) \right].
\end{equation}

\begin{equation}
\label{eq:hjb-ell}
    \frac{\partial v}{\partial t} + \frac{\w{\sigma}^2(t)}{2} \Delta v +\w{b} (t,x) \cdot \nabla v + \inf_{u } \left\{ \ell(t,x,u) + u \cdot \nabla v \right\} = 0, \quad v(T,x) = r(x).
\end{equation}

\begin{theorem}[Regularity and Well-Posedness]
\label{theo:Regularity}
Suppose that Assumptions~\textnormal{(H1)--(H7)} hold.  
Then the generalized fine-tuning (GFT) problem~\ref{prob:GFT} is well-posed, and the following properties hold:

\begin{itemize}
\item[\textbf{(i)}]  
There exists a unique classical solution 
\[
v^*(t,x) \in C^{1,2}([0,T]\times\mathbb{R}^d)
\]
to the Hamilton--Jacobi--Bellman (HJB) equation~\eqref{eq:hjb-ell}.  
This function coincides with the optimal value function defined in~\eqref{eq:v^*}.  
Moreover, there exist constants $L_1,L_2<\infty$, depending only on the structural parameters $(\lambda_\ell,L_b,L_\sigma,L_\ell,L_r)$, such that
\[
\|\nabla_x v^*\|_\infty \le L_1, 
\qquad 
\|\nabla_{xx}^2 v^*\|_\infty \le L_2,
\]
and consequently $\|v^*(t,\cdot)\|_{\mathrm{Lip}}\le L_1$ for all $t\in[0,T]$.

\item[\textbf{(ii)}] 
There exists a unique optimal feedback control 
\[
u^*(t,x) = \arg\min_{u\in\mathbb{R}^d}\{\ell(t,x,u)+u\cdot\nabla_x v^*(t,x)\}
\]
solving Problem~\ref{prob:GFT}.  
The feedback map $u^*(t,x)$ is continuously differentiable in $x$ and globally Lipschitz; in particular,
\[
u^* \in C^{0,1}([0,T]\times\mathbb{R}^d;\mathbb{R}^d),
\qquad 
\nabla_x u^* \in L^\infty([0,T]\times\mathbb{R}^d),
\]
and $\nabla_x u^*(t,x)$ is Lipschitz continuous in $x$ uniformly over $t$.
\end{itemize}
\end{theorem}

\begin{proof}[Proof of Theorem~\ref{theo:Regularity}]
We work under Assumptions~(H1)--(H7).  
Fix $\beta>0$ and define the weighted Lipschitz norm
\[
\|u\|_{\mathrm{Lip},\beta}
=\sup_{t\in[0,T]}e^{-\beta t}\sup_{x\neq y}\frac{\|u(t,x)-u(t,y)\|}{\|x-y\|}.
\]
Let $\mathcal{B}_R:=\{u:\|u\|_{\mathrm{Lip},\beta}\le R\}$.

\paragraph{Step 1. (Regularity for fixed $u$).}
For $u\in\mathcal{B}_R$, consider the FBSDE
\[
\begin{cases}
dX_s^{t,x}=(\w b(s,X_s^{t,x})+u(s,X_s^{t,x}))\,ds+\w\sigma(s)\,dB_s,\\[3pt]
dY_s^{t,x}=-\ell(s,X_s^{t,x},u(s,X_s^{t,x}))\,ds+Z_s^{t,x}\,dB_s,\\[3pt]
Y_T^{t,x}=r(X_T^{t,x}).
\end{cases}
\]
Under (H1)--(H7), all coefficients are $C^2$ in $x$, globally Lipschitz, and $\w\sigma$ is uniformly elliptic.  
By the Four-Step Scheme \textnormal{(Ma, Yong \& Zhang 1999, Thm.~5.3.3)} or the regularity theory for coupled FBSDEs \textnormal{(Delarue 2002, Thm.~2.1)}, there exists a unique decoupling field $v^u\in C^{1,2}$ satisfying
\[
\partial_t v^u + (\w b+u)\cdot\nabla_x v^u
+ \tfrac{1}{2}\operatorname{Tr}(\w\sigma\w\sigma^\top\nabla_{xx}^2 v^u)
+ \ell(t,x,u)=0,\qquad v^u(T,x)=r(x),
\]
together with uniform bounds (over $(t,x)\in[0,T]\times\mathbb{R}^d$):
\[
\sup_{t,x}\|\nabla_x v^u(t,x)\|\le C_1(R),\qquad
\sup_{t,x}\|\nabla_{xx}^2 v^u(t,x)\|\le C_2(R),
\]
for some continuous functions $C_1,C_2$ depending on $R=\|u\|_{\mathrm{Lip},\beta}$.

\paragraph{Step 2. (Definition of $\Phi$ and self-mapping).}
Given $v^u$, define
\[
\nabla_u\ell(t,x,u'(t,x))+\nabla_x v^u(t,x)=0,\qquad
u'=\Phi(u).
\]
By strong convexity of $\ell$ in $u$ ($\lambda_\ell>0$), there exists a unique minimizer $u'(t,x)$ for each $(t,x)$ (\emph{pointwise uniqueness}).  
By the implicit function theorem (since $\nabla_{uu}^2\ell$ is invertible with eigenvalues $\ge\lambda_\ell$), $u'\in C^1$ in $x$.  
Differentiating yields
\[
\nabla_x u'(t,x)
=-(\nabla_{uu}^2\ell)^{-1}(\nabla_{xu}^2\ell+\nabla_{xx}^2 v^u),
\]
so that
\[
\|\nabla_x u'\|_\infty
\le\tfrac{1}{\lambda_\ell}(L_{xu}^\ell+C_2(R))
+\tfrac{L_{uu}^\ell}{\lambda_\ell^2}(L_{xu}^\ell+C_2(R)).
\]
Choosing 
\[
R \ge R_0 := 
\tfrac{1}{\lambda_\ell}(L_{xu}^\ell + C_2(R)) 
+ \tfrac{L_{uu}^\ell}{\lambda_\ell^2}(L_{xu}^\ell + C_2(R))
\]
ensures $\Phi(\mathcal{B}_R)\subset\mathcal{B}_R$.

\paragraph{Lemma 1 (Forward stability).}
Let $u^1,u^2\in\mathcal{B}_R$.  
If $X^i$ solve their respective forward SDEs, then
\[
\sup_{s\in[t,T]}\mathbb{E}\|X_s^1-X_s^2\|^2
\le C_\beta\,\|u^1-u^2\|_{\mathrm{Lip},\beta}^2,
\qquad
C_\beta=\tfrac{C(e^{2\beta T}-1)}{\beta}.
\]

\emph{Proof.}  
By Itô’s formula,
\[
d\|X^1_s-X^2_s\|^2
\le C\|X^1_s-X^2_s\|^2\,ds + C\|u^1(s,X^1_s)-u^2(s,X^2_s)\|^2\,ds + dM_s.
\]
Using the triangle inequality and Lipschitz continuity,
\[
\|u^1(s,X^1_s)-u^2(s,X^2_s)\|
\le \|u^1(s,X^1_s)-u^1(s,X^2_s)\| + \|u^1(s,X^2_s)-u^2(s,X^2_s)\|
\le e^{\beta s}\|u^1-u^2\|_{\mathrm{Lip},\beta}(\|X^1_s-X^2_s\|+\|X^2_s\|).
\]
Taking expectations, using the linear growth bound 
$\mathbb{E}\|X^2_s\|^2 \le C(1+|x|^2)$, 
and applying Gronwall’s inequality yields the result. $\square$

\paragraph{Lemma 2 (BSDE stability).}
Let $(Y^i,Z^i)$ correspond to $u^i$. Then
\[
\sup_{s\in[t,T]}\mathbb{E}\|Y_s^1-Y_s^2\|^2
+ \mathbb{E}\!\int_t^T\!\|Z_s^1-Z_s^2\|^2\,ds
\le C_\beta'\,\|u^1-u^2\|_{\mathrm{Lip},\beta}^2,
\]
where $C_\beta'$ depends on $C_\beta$ from Lemma~1.
\emph{Proof.}  
Subtract the BSDEs and apply standard $L^2$ energy estimates (\textnormal{Pardoux \& Peng 1990, Lemma~2.1}) and Lemma~1. $\square$

\paragraph{Lemma 3 (Gradient stability).}
Let $Y^{i,x}_s:=v^i(s,X^{i,t,x}_s)$, where $X^{i,t,x}$ starts from $(t,x)$.  
Define the difference quotient 
\[
\Delta^\varepsilon Y^i_s := \tfrac{Y^{i,x+\varepsilon h}_s-Y^{i,x}_s}{\varepsilon},\qquad |h|=1.
\]
Then $(\Delta^\varepsilon Y^i_s,\Delta^\varepsilon Z^i_s)$ satisfies a linear variational BSDE obtained by differentiating the original BSDE w.r.t.\ the initial condition~$x$ (see \textnormal{Delarue 2002, Lemma~2.4}).  
By BSDE stability,
\[
E\sup_{s\in[t,T]}|\Delta^\varepsilon Y^1_s-\Delta^\varepsilon Y^2_s|^2
\le\tfrac{C(e^{2\beta T}-1)}{\beta}\|u^1-u^2\|_{\mathrm{Lip},\beta}^2.
\]
Letting $\varepsilon\downarrow0$, by the $C^{1,2}$ regularity of $v^i$ (Step~A0), we have 
$\Delta^\varepsilon Y^i_s \to \nabla_x v^i(s,X^{i,t,x}_s)\cdot h$ in $L^2$.  
Using the martingale representation $Z_s^i=\w\sigma(s)^\top\nabla_x v^i(s,X^i_s)$ gives
\[
\sup_{s,x}\|\nabla_x v^1(s,x)-\nabla_x v^2(s,x)\|
\le K(\beta)\|u^1-u^2\|_{\mathrm{Lip},\beta},\qquad K(\beta)\le\tfrac{C}{\beta}.
\]
The factor $1/\beta$ arises from the time integral 
$\int_s^T e^{\beta r}\,dr=(e^{\beta T}-e^{\beta s})/\beta$ in the weighted norm. $\square$

\paragraph{Step 3. (Contraction and uniqueness).}
For $u'^i=\Phi(u^i)$,
\[
\nabla_u\ell(t,x,u'^1)-\nabla_u\ell(t,x,u'^2)
=-(\nabla_x v^1-\nabla_x v^2).
\]
By strong monotonicity of $\nabla_u\ell$,
\[
\|u'^1-u'^2\|
\le\tfrac{1}{\lambda_\ell}\|\nabla_x v^1-\nabla_x v^2\|,
\]
thus
\[
\|\Phi(u^1)-\Phi(u^2)\|_{\mathrm{Lip},\beta}
\le\tfrac{K(\beta)}{\lambda_\ell}\|u^1-u^2\|_{\mathrm{Lip},\beta}.
\]
Since $K(\beta)\le C/\beta$, choosing $\beta>C/\lambda_\ell$ makes $\Phi$ a strict contraction.  
Hence a unique fixed point $u^*\in\mathcal{B}_R$ exists (\emph{functional uniqueness}),  
and the associated decoupling field $v^*=v^{u^*}\in C^{1,2}$ inherits the bounded derivative properties of Step~A0.

\paragraph{Conclusion.}
The fixed-point construction establishes Theorem~\ref{theo:Regularity}:  
(i) uniqueness and regularity of $v^*$ as in part~(i) of the theorem;  
(ii) uniqueness and regularity of $u^*$ as in part~(ii);  
and (iii) all constants depend only on the structural parameters, not on $(t,x)$.
\end{proof}

\end{document}
